% Initial research-talk deck based on overleaf/main_linear.tex.
% Remove `handout` to enable overlays when presenting.
\documentclass[handout,compress,9pt]{beamer}

\usefonttheme[onlymath]{serif}
\usepackage[english]{babel}
\usepackage{amsmath,amssymb,amsthm}
\usepackage{bm}
\usepackage{booktabs}
\usepackage{graphicx}
\usepackage{tikz}
\usetikzlibrary{arrows.meta,positioning,calc}

\usetheme{Boadilla}
\useoutertheme{infolines}
\setbeamertemplate{navigation symbols}{}
\setbeamertemplate{theorems}[numbered]
\setbeamertemplate{caption}[numbered]

\definecolor{ddlblue}{RGB}{0,78,140}
\definecolor{ddlred}{RGB}{176,34,34}
\definecolor{ddlgreen}{RGB}{32,120,82}
\newcommand{\blue}[1]{\textcolor{ddlblue}{#1}}
\newcommand{\red}[1]{\textcolor{ddlred}{#1}}
\newcommand{\green}[1]{\textcolor{ddlgreen}{#1}}

\newcommand{\RR}{\mathbb{R}}
\newcommand{\zero}{\bm{0}}
\newcommand{\bx}{\bm{x}}
\newcommand{\by}{\bm{y}}
\newcommand{\bz}{\bm{z}}
\newcommand{\bw}{\bm{w}}
\newcommand{\br}{\bm{r}}
\newcommand{\bmu}{\bm{\mu}}
\newcommand{\blmbd}{\bm{\lambda}}
\newcommand{\calA}{\mathcal{A}}
\newcommand{\calL}{\mathcal{L}}
\newcommand{\calO}{\mathcal{O}}
\newcommand{\calS}{\mathcal{S}}
\newcommand{\prox}{\operatorname{prox}}
\newcommand{\dist}{\operatorname{dist}}
\newcommand{\st}{\textnormal{s.t.}}
\newcommand{\eps}{\varepsilon}

\setbeamertemplate{footline}{%
  \leavevmode%
  \hbox{%
    \begin{beamercolorbox}[wd=.23\paperwidth,ht=2.25ex,dp=1ex,center]{author in head/foot}%
      \usebeamerfont{author in head/foot}\insertshortauthor~(\insertshortinstitute)
    \end{beamercolorbox}%
    \begin{beamercolorbox}[wd=.54\paperwidth,ht=2.25ex,dp=1ex,center]{title in head/foot}%
      \usebeamerfont{title in head/foot}\insertshorttitle
    \end{beamercolorbox}%
    \begin{beamercolorbox}[wd=.23\paperwidth,ht=2.25ex,dp=1ex,right]{date in head/foot}%
      \usebeamerfont{date in head/foot}\insertshortdate{}\hspace*{1.4em}\insertframenumber{} / \inserttotalframenumber\hspace*{1ex}
    \end{beamercolorbox}%
  }%
  \vskip0pt%
}

\title[Dual-Damping Linearized ALM]{Dual-Damping Linearized ALM\\for Nonsmooth Composite Optimization}
\author[L. Li et al.]{Luxuan Li\inst{1}}
%  \and Kangkang Deng\inst{2} \and Chunfeng Cui\inst{1} \and Jiang Hu\inst{3}}
\institute[Beihang]{\inst{1} Beihang University}
%  \and \inst{2} National University of Defense Technology \and \inst{3} Tsinghua University}
\date[Research update]{Research update, August 2026}

% Insert a standalone divider slide before every section.
\AtBeginSection[]{
  \begin{frame}[plain,noframenumbering]
    \centering
    \vfill
    {\usebeamerfont{section title}\usebeamercolor[fg]{section title}%
      \Roman{section}.\enspace\insertsectionhead\par}
    \vfill
  \end{frame}
}

\begin{document}

\begin{frame}
  \titlepage
\end{frame}

\begin{frame}{Outline}
  \tableofcontents
\end{frame}

\section{Problem Setting}

\begin{frame}{Structured Nonsmooth Composite Optimization}
We consider
\begin{equation*}
  \min_{\bx\in\RR^n}\quad F(\bx):=f(\bx)+h\bigl(\calA(\bx)\bigr).
  \tag{P0}
\end{equation*}
\begin{itemize}
  \item $f:\RR^n\to\RR$ is continuously differentiable and may be nonconvex.
  \item $\calA:\RR^n\to\RR^m$ is continuously differentiable and may be nonlinear.
  \item $h:\RR^m\to\RR$ is convex, Lipschitz continuous, and possibly nonsmooth.
\end{itemize}
\end{frame}

\begin{frame}{Approximate KKT Point}
  \begin{itemize}
  \item Introduce $\by=\calA(\bx)$:
  \begin{equation}
    \label{Prob:split}
    \min_{\bx,\by}\quad f(\bx)+h(\by)
    \qquad \st\qquad \calA(\bx)-\by=\zero.   \tag{P}
  \end{equation}

  \item The split KKT system is
  \begin{equation*}
    \nabla f(\bx)+J_{\calA}(\bx)^\top\bz=\zero,
    \qquad \bz\in\partial h(\by),
    \qquad \calA(\bx)-\by=\zero.
  \end{equation*}

  \item \textbf{approximate $\eps$-KKT point.}
  Find $(\bx,\by,\bz)$ such that
  \begin{equation*}
    \|\nabla f(\bx)+J_{\calA}(\bx)^\top\bz\|\leq\eps,
    \quad \blue{\bz\in\partial h(\by)},
    \quad \|\calA(\bx)-\by\|\leq\eps.
  \end{equation*}

  \item The exact $\by$-proximal step makes the subdifferential condition \blue{exact}.
  \end{itemize}
\end{frame}

\begin{frame}{Main Results}
\centering
{\scriptsize
\renewcommand{\arraystretch}{1.35}
\setlength{\tabcolsep}{1.2pt}
\begin{tabular}{@{}p{0.14\textwidth}p{0.12\textwidth}p{0.07\textwidth}p{0.11\textwidth}p{0.055\textwidth}p{0.085\textwidth}p{0.075\textwidth}p{0.19\textwidth}p{0.09\textwidth}@{}}
\toprule
\textbf{Method}
& \shortstack{\textbf{Constrained}\\\textbf{setting}}
& $A$
& \shortstack{\textbf{Nonsmoothness}\\\textbf{of $h$}}
& \textbf{Loop}
& \shortstack{KKT\\measure}
& \shortstack{Penalty\\$\rho$}
& \shortstack{Dual regularizer\\ $s$}
& \shortstack{\textbf{Iteration}\\\textbf{complexity}}
\\
\midrule
MRSALM
& Manifold
& Linear
& Lipschitz
& $\checkmark$
& $\eps$-KKT
& fixed
& $s_k=\Theta(k^{1/4})\uparrow$
& $O(\eps^{-4})$
\\
\shortstack{MRALM}
& Manifold
& Linear
& Lipschitz
& $\times$
& $\eps$-KKT
& fixed
& $s_k=\Theta(k^{1/3})\uparrow$
& $O(\eps^{-3})$
\\
\blue{DDLALM (Ours)}
& \blue{\shortstack{Uncon-\\strained}}
& \blue{\shortstack{Non-\\linear}}
& \blue{Lipschitz}
& \blue{$\checkmark$}
& \blue{$\eps$-KKT}
& fixed
& \blue{$s_k=\Theta(k^{1/3})\uparrow$}
& \blue{$O(\eps^{-3})$}
\\
\bottomrule
\end{tabular}}
\end{frame}

\section{Dual-Damping Reformulation}

\begin{frame}{Dual-Proximal Augmented Lagrangian}
\begin{itemize}
  \item The Augmented Lagrangian of Problem \eqref{Prob:split} is defined as
  \[
  \calL_\rho(\bx,\by,\blmbd)
=f(\bx)+h(\by)+\langle\blmbd,\calA(\bx)-\by\rangle
+\frac{\rho}{2}\|\calA(\bx)-\by\|^2.
  \]

\item \blue{Dual-regularized Augmented Lagrangian}: Fix a center multiplier $\bmu$ and introduce $s>0$, we define
\begin{equation*}
  \mathcal P_s(\bx,\by,\blmbd;\bmu)
  :=\calL_\rho(\bx,\by,\blmbd)
  -\blue{\frac{1}{2s}\|\blmbd-\bmu\|^2}.
\end{equation*}

\item For fixed $(\bx,\by,\bmu)$, $\blmbd\mapsto\mathcal P_s$ is $1/s$-strongly concave,
  admitting a unique maximizer
\begin{equation*}
  \blue{\blmbd_s(\bx,\by;\bmu)=\bmu+s\bigl(\calA(\bx)-\by\bigr).}
\end{equation*}

\item nonconvex--concave $\min_{x,y}\max_\lambda \mathcal L_\rho(x,y,\lambda)$ $\longrightarrow$

nonconvex--\blue{strongly concave} problem:
\begin{equation*}
  \min_{\bx,\by}\max_{\blmbd}
  \mathcal P_s(\bx,\by,\blmbd;\bmu).
\end{equation*}

\item  \textbf{Reduced Augmented Lagrangian.}
Substitution $\blmbd_s(\bx,\by;\bmu)$ into $\mathcal P_s$ yields the value function:
\begin{equation*}
  \begin{aligned}
  \Phi_s(\bx,\by;\bmu)
  &:= \max_{\blmbd}\mathcal P_s(\bx,\by,\blmbd;\bmu) \\
  &=f(\bx)+h(\by)+\langle\red{\bmu},\calA(\bx)-\by\rangle
  +\frac{\rho+s}{2}\|\calA(\bx)-\by\|^2.
  \end{aligned}
\end{equation*}
\end{itemize}
\end{frame}

\begin{frame}{Danskin's Theorem for the Value Function $ \Phi_s(\bx,\by;\bmu)$}
% \scriptsize
% The value function:
% \begin{equation*}
%   \begin{aligned}
%   \Phi_s(\bx,\by;\bmu)
%   &:= \max_{\blmbd}\mathcal P_s(\bx,\by,\blmbd;\bmu) \\
%   &=f(\bx)+h(\by)+\langle\red{\bmu},\calA(\bx)-\by\rangle
%   +\frac{\rho+s}{2}\|\calA(\bx)-\by\|^2.
%   \end{aligned}
% \end{equation*} with 
% Strong concavity in $\blmbd$ gives the unique maximizer
% \begin{equation*}
%   \blue{\blmbd_s(\bx,\by;\bmu)
%   =\bmu+s\bigl(\calA(\bx)-\by\bigr).}
% \end{equation*}
Danskin's theorem gives the smooth outer gradients: \hyperlink{appendix:danskin-theorem}{%
  \beamergotobutton{Theorem $\rightarrow$}%
}
\begin{itemize}
\setlength{\itemsep}{0.10cm}

\item[\blue{$\triangleright$}]
\textbf{$\bx$-gradient:}
\begin{equation*}
\begin{aligned}
  \nabla_{\bx}\Phi_s(\bx,\by;\bmu)
  &=\nabla_{\bx}\mathcal P_s
  \bigl(\bx,\by,\blmbd_s(\bx,\by;\bmu);\bmu\bigr)\\
  &=\nabla f(\bx)+J_{\calA}(\bx)^\top
  \left[\bmu+(\rho+s)\bigl(\calA(\bx)-\by\bigr)\right].
\end{aligned}
\end{equation*}

\item[\blue{$\triangleright$}]
\textbf{$\by$-subgradient:}
\begin{equation*}
  \partial_{\by}\Phi_s(\bx,\by;\bmu)
  =\partial h(\by)-\bmu-(\rho+s)\bigl(\calA(\bx)-\by\bigr).
\end{equation*}

\item[\blue{$\triangleright$}]
\textbf{$\bmu$-gradient:}
\begin{equation*}
\begin{aligned}
  \nabla_{\bmu}\Phi_s(\bx,\by;\bmu)
  &=\nabla_{\bmu}\mathcal P_s
  \bigl(\bx,\by,\blmbd_s(\bx,\by;\bmu);\bmu\bigr)\\
  &=\frac{1}{s}\bigl(\blmbd_s(\bx,\by;\bmu)-\bmu\bigr)
  =\calA(\bx)-\by.
\end{aligned}
\end{equation*}
\end{itemize}
\end{frame}

\section{Dual-Damping Linearized ALM}

\begin{frame}{Core Framework: Coupled Saddle Update}
\vspace{0.1cm}

\noindent\textbf{Idea:} Apply gradient steps on
$\Phi_{s_k}$ and solve the coupled $(\by,\blmbd)$ saddle block exactly.

\noindent\textbf{Parameter:} tolerance $\varepsilon > 0 $, $\eta_k, \, \tau_k, \, s_k$, max iteration number $K$.

Set $\br^k:=\calA(\bx^k)-\by^k$ and $\rho_{s_k}:=\rho+s_k$.

\vspace{0.12cm}
\begin{itemize}
\item[\blue{$\triangleright$}] \textbf{Steps 1--2: $\bx$-gradient descent.}
\begin{equation}\label{eq:slide-exact-multiplier-x}
\begin{aligned}
  % \blmbd_{s_k}
  % &=\arg\max_{\blmbd}\mathcal P_{s_k}
  % (\bx^k,\by^k,\blmbd;\bmu^k)
  % =\bmu^k+s_k\br^k,\\
  \bx^{k+1}
  &=\bx^k-\eta_k\nabla_{\bx}
  \Phi_{s_k}(\bx^k,\by^k;\bmu^k)\\
  &=\bx^k-\eta_k\left[
  \nabla f(\bx^k)+J_{\calA}(\bx^k)^\top
  \bigl(\bmu^k+\rho_{s_k}\br^k\bigr)\right].
\end{aligned}
\end{equation}

\vspace{0.1cm}
\item[\blue{$\triangleright$}] \textbf{Step 3: coupled $(\by,\blmbd)$ saddle update.}
\begin{equation}\label{eq:slide-coupled-saddle-update}
  (\by^{k+1},\blmbd^{k+1})
  =
  \operatorname{saddle}_{\by,\blmbd}
  \mathcal P_{s_k}(\bx^{k+1},\by,\blmbd;\bmu^k).
\end{equation}

\vspace{0.1cm}
\item[\blue{$\triangleright$}] \textbf{Step 4: $\bmu$-gradient ascent on $\Phi_{s_k}$.}
\begin{equation}\label{eq:slide-mu-update}
\begin{aligned}
  \bmu^{k+1}
  &=\bmu^k+\tau_k\nabla_{\bmu}
  \Phi_{s_k}(\bx^{k+1},\by^{k+1};\bmu^k)\\
  &=\bmu^k+\tau_k\br^{k+1}.
\end{aligned}
\end{equation}
\end{itemize}
\end{frame}

\begin{frame}{Step 3: Coupled $(\by,\blmbd)$ Saddle Update}
\vspace{0.2cm}
With $(\bx^{k+1},\bmu^k)$ fixed, solve the saddle subproblem exactly:
\begin{equation*}
  \min_{\by}\max_{\blmbd}\;
  \mathcal P_{s_k}(\bx^{k+1},\by,\blmbd;\bmu^k).
\end{equation*}

\vspace{0.1cm}
\noindent\textbf{Step 3: Executable form:}\hyperlink{appendix:why-rho-plus-sk}{%
  \beamergotobutton{Why $\rho+s_k$?}%
}

\begin{equation*}
  \begin{aligned}
  \by^{k+1}
  &=
  \prox_{h/\rho_{s_k}}
  \left(\calA(\bx^{k+1})+\frac{\bmu^k}{\rho_{s_k}}\right), \\
  \blmbd^{k+1}
  &=
  \bmu^k+s_k\br^{k+1}.
  \end{aligned}
\end{equation*}

\begin{itemize}

\vspace{0.06cm}
\item[\blue{$\triangleright$}] \textbf{Coupled saddle point:}
\begin{equation*}
  \begin{aligned}
  \by^{k+1}
  &=\operatorname*{arg\,min}_{\by}\;
  \mathcal P_{s_k}(\bx^{k+1},\by,\blmbd^{k+1};\bmu^k) =
  \prox_{h/\rho}
  \left(\calA(\bx^{k+1})+\frac{\red{\blmbd^{k+1}}}{\rho}\right),
  \\
  \red{\blmbd^{k+1}}
  &=\operatorname*{arg\,max}_{\blmbd}\;
  \mathcal P_{s_k}(\bx^{k+1},\by^{k+1},\blmbd;\bmu^k) =  \bmu^k+s_k\bigl(\calA(\bx^{k+1})-\by^{k+1}\bigr).
  \end{aligned}
\end{equation*}

\vspace{0.06cm}
\item[\blue{$\triangleright$}] \textbf{Optimality conditions:} 
\begin{equation*}
  0\in
  \partial h(\by^{k+1})
  -\blmbd^{k+1}
  +\rho\bigl(\by^{k+1}-\calA(\bx^{k+1})\bigr),
\end{equation*}
\begin{equation*}
  0=
  \calA(\bx^{k+1})-\by^{k+1}
  -\frac{1}{s_k}\bigl(\blmbd^{k+1}-\bmu^k\bigr).
\end{equation*}

\end{itemize}
\end{frame}

\begin{frame}{Key Identities Used Throughout}

\begin{itemize}
  \item[\blue{$\triangleright$}]\textbf{Exact identities}:
\begin{equation}\label{eq:key-exact-identities}
\begin{aligned}
  \blmbd^{k+1}
  &=\bmu^k+s_k\br^{k+1},\\
  \bmu^{k+1}-\bmu^k
  &=\tau_k\br^{k+1}
  =\frac{\tau_k}{s_k}(\blmbd^{k+1}-\bmu^k),\\
  \blmbd^{k+1}+\rho\br^{k+1}
  &=\bmu^k+\rho_{s_k}\br^{k+1}
  \in\partial h(\by^{k+1}).
\end{aligned}
\end{equation}

\item[\blue{$\triangleright$}]\textbf{Convex-combination identities}
\begin{equation}\label{eq:key-convex-combinations}
\begin{aligned}
  \blmbd^{k+1}
  &=\frac{\rho}{\rho_{s_k}}\bmu^k
  +\frac{s_k}{\rho_{s_k}}
  (\blmbd^{k+1}+\rho\br^{k+1}),\\
  \bmu^{k+1}
  &=\left(1-\frac{\tau_k}{s_k}\right)\bmu^k+ 
  \frac{\tau_k}{s_k}\blmbd^{k+1}.
\end{aligned}
\end{equation}
\end{itemize}

\end{frame}

\section{Complexity Analysis}

\begin{frame}{Convergence: $O(\eps^{-3})$ Iteration Complexity}
Under mild assumptions, choose
\begin{equation*}
  s_k=c_s(k+1)^{1/3},
  \quad \eta_k=\frac{c_x}{s_k},
  \quad \tau_k=\frac{\theta_\mu s_k}{k+2} = \Theta(k^{-2/3}).
  \quad c_x\widehat L\left(1+\frac{\rho}{s_0}\right)\le \frac12.
\end{equation*}
where $c_s$, $c_x$, $\theta_\mu$, $\rho$, and $\widehat L$ are positive constants independent of $k$.
Then there exists an index
$R\in\{\lceil K/2\rceil,\ldots,K-1\}$ such that
$\blmbd^R+\rho\br^R\in\partial h(\by^R)$ and
\begin{equation*}
  \left\|\nabla f(\bx^R)+J_{\calA}(\bx^R)^\top
  \bigl(\blmbd^R+\rho\br^R\bigr)\right\|
  =O\bigl(K^{-1/3}\bigr),
  \qquad
  \|\calA(\bx^R)-\by^R\|=O(K^{-1/3}).
\end{equation*}
Consequently, the oracle complexity of DDLALM to reach an approximate $\varepsilon$-KKT is of the order $\calO(\varepsilon^{-3})$.
\end{frame}

\begin{frame}[t]{Main Assumptions}
\small
\vspace{0.15cm}

% \noindent\textbf{Assumptions}
\begin{itemize}
\item[(i)] \textbf{Level-bounded objective.}
For every $\alpha\in\RR$, the set
\begin{equation*}
  \calS_{\alpha}:=\{\bx\in\RR^n:F(\bx)\leq\alpha\}
\end{equation*}
is empty or compact.

\item[(ii)] \textbf{Smoothness and boundedness\footnote{Y. Xie and S. J. Wright, ``Complexity of Proximal Augmented Lagrangian for Nonconvex Optimization with Nonlinear Equality Constraints,'' \emph{J. Sci. Comput.}, 86:38, 2021.}.}
Given a nonempty working level set $\calS:=\calS_{\alpha_0}$.
The functions $f$ and $\calA$ are $L_f$- and $L_{\calA}$-smooth on $\calS$.
Moreover,
\begin{equation*}
  \|\nabla f(\bx)\|\leq G_f,
  \quad \|J_{\calA}(\bx)\|\leq G_{\calA},
  \quad M_{\calA}:=\sup_{\bx\in\calS}\|\calA(\bx)\|<\infty.
\end{equation*}

\item[(iii)] $h$ is proper, closed, convex, and globally $L_h$-Lipschitz; equivalently,
\begin{equation*}
  \|\bz\|\leq L_h
  \qquad \text{for all }\bz\in\partial h(\by).
\end{equation*}
\end{itemize}

\vspace{0.15cm}
\begin{block}{Remark}
The smoothness condition can be imposed on a compact convex set $\calS$
satisfying $\calS_\alpha\subseteq\operatorname{int}(\calS)$.
The containment $\bx^k\in\calS_\alpha$ for every $k\ge0$ is
\red{established by induction}.
\hfill
\hyperlink{appendix:bootstrap-joint-claim}{%
  \beamergotobutton{Joint induction claim $\rightarrow$}%
}
\end{block}
\end{frame}

\begin{frame}{Proof Roadmap}
\scriptsize
\vspace{0.1cm}
\begin{itemize}
\setlength{\itemsep}{0.20cm}

\item[\red{$\blacktriangleright$}]
\textbf{Step 1: establish multiplier boundedness and feasibility control.}

Assumption (iii) provides bounded subgradients of $h$.
The two convex-combination identities in
\eqref{eq:key-convex-combinations} imply
\begin{equation*}
  \|\blmbd^{k+1}\|,\ \|\bmu^{k+1}\|\le M_\mu,
  \qquad
  \|\br^{k+1}\|\le \frac{C_r}{\rho_{s_k}}.
\end{equation*}

\item[\red{$\blacktriangleright$}]
\textbf{Step 2: sufficient primal descent on the value function $\Phi_s$.}

The $\bx$-gradient step and the exact $(\by,\blmbd)$ saddle solve in in \eqref{eq:slide-coupled-saddle-update} yield
sufficient primal descent; \red{Exactness eliminates the multiplier-tracking error with respect to $\blmbd$.} Together, these properties yield
\begin{equation*}
  \Phi_{s_k}(\bx^{k+1},\by^{k+1};\bmu^k)
  -\Phi_{s_k}(\bx^k,\by^k;\bmu^k)
  \le -\frac{3\rho_{s_k}}{4\theta_x}\|\Delta_{\bx}^{k+1}\|^2
  -\frac{\rho_{s_k}}{2}\|\Delta_{\by}^{k+1}\|^2.
\end{equation*}
\vspace{-0.22cm}

\item[\red{$\blacktriangleright$}]
\textbf{Step 3: control the moving-potential drift.}
Using $\Delta_{\bmu}^{k+1}=\tau_k\br^{k+1}$ from
\eqref{eq:key-exact-identities}, the two shifts ($s_k$ and $\bmu^k$) satisfy
\vspace{-0.08cm}
\begin{equation*}
\begin{aligned}
  &\Phi_{s_{k+1}}(\bx^{k+1},\by^{k+1};\bmu^{k+1})
  -\Phi_{s_k}(\bx^{k+1},\by^{k+1};\bmu^k) = 
  \underbrace{\langle\Delta_{\bmu}^{k+1},\br^{k+1}\rangle
  +\frac{\rho_{s_{k+1}}-\rho_{s_k}}{2}\|\br^{k+1}\|^2}_{\text{moving-potential drift}}\\
  &\qquad\le
  \underbrace{C_r^2\!\left(
  \frac{\tau_k}{\rho_{s_k}^2}
  +\frac{\rho_{s_{k+1}}-\rho_{s_k}}{2\rho_{s_k}^2}
  \right)}_{e_k}
\end{aligned}
\end{equation*}
\vspace{-0.08cm}

\item[\red{$\blacktriangleright$}]
\textbf{Step 4: choose the parameter schedules and control the KKT residuals.}

The parameter schedules make the drift errors summable:
\begin{equation*}
  \sum_{k=0}^{\infty} e_k < \infty
\end{equation*}
\vspace{-0.18cm}
\hfill
\hyperlink{appendix:lemmas}{%
  \beamergotobutton{Appendix I: Lemmas 1--4 $\rightarrow$}%
}

\end{itemize}


\end{frame}

% \begin{frame}{Theorem: An $O(1/K)$ Averaged Bound}
% \small
% Lemma 4 gives the finite weighted budget
% \begin{equation*}
%   c_{\mathrm{desc}}
%   \sum_{k=0}^{K-1}\rho_{s_k}D_k
%   \le C_{\mathrm{desc}},
%   \qquad
%   D_k:=\|\Delta_{\bx}^{k+1}\|^2
%   +\|\Delta_{\by}^{k+1}\|^2.
% \end{equation*}

% Dividing both sides by $K$ gives
% \begin{block}{Averaged rate}
% \begin{equation*}
%   \blue{\boxed{
%   \frac{1}{K}\sum_{k=0}^{K-1}\rho_{s_k}D_k
%   \le \frac{C_{\mathrm{desc}}}{c_{\mathrm{desc}}K}
%   =O\!\left(\frac{1}{K}\right)}}.
% \end{equation*}
% \end{block}

% \end{frame}

\begin{frame}{Comparison with Closely Related Methods}
\centering
{\scriptsize
\renewcommand{\arraystretch}{1.35}
\setlength{\tabcolsep}{1.2pt}
\resizebox{\textwidth}{!}{%
\begin{tabular}{@{}p{0.14\textwidth}p{0.12\textwidth}p{0.11\textwidth}p{0.105\textwidth}p{0.10\textwidth}p{0.17\textwidth}p{0.09\textwidth}@{}}
\toprule
\textbf{Method}
& \shortstack{\textbf{Problem}\\\textbf{setting}}
& \shortstack{\textbf{Convexity}\\\textbf{of $h$}}
& \shortstack{\textbf{KKT}\\\textbf{measure}}
& \shortstack{\textbf{Penalty}\\$\boldsymbol{\rho}$}
& \shortstack{\textbf{Parameter}\\\textbf{schedule}}
& \shortstack{\textbf{Iteration}\\\textbf{complexity}}
\\
\midrule
\shortstack{SDCAM$_{1\ell}$\\\green{[1]}}
& Composite
& $\checkmark$
& \shortstack{$(\eps_1,\eps_2,0)$\\stationarity}
& N/A
& \shortstack{Continuation\\$\uparrow$}
& $O(\eps^{-3})$
\\
\shortstack{SDD-ADMM\\\green{[2]}}
& \shortstack{Nonlinear\\equality}
& $\checkmark$(ncvx)
& \shortstack{$\eps$-\\stationarity}
& $\Theta(\eps^{-1})$
& \shortstack{$\omega,\tau$ fixed\\per run}
& \shortstack{$O(\eps^{-3})$}
\\
\blue{\shortstack{\textbf{DDLALM}\\\textbf{[This work]}}}
& \blue{$h(\calA(\bx))$}
& \blue{$\checkmark$}
& \blue{\shortstack{\textbf{split}\\$\boldsymbol{\eps}$\textbf{-KKT}}}
& \blue{\textbf{fixed}}
& \blue{\shortstack{$s_k=\Theta(k^{1/3})$\\$\uparrow$}}
& \blue{\textbf{$O(\eps^{-3})$}}
\\
\bottomrule
\end{tabular}%
}}

\vspace{2mm}
{\tiny
SDCAM$_{1\ell}$ is the closest same-scope complexity comparator;
SDD-ADMM is the nonlinear ALM/ADMM-family comparator.
}
\end{frame}

\section{Discussion}

\begin{frame}{Concern I: Can the $\by$-Update Avoid $s_k$ in DDALM?}
Recalling
\begin{equation*}
  \begin{aligned}
  \mathcal P_{s_k}(\bx,\by,\blmbd;\bmu)
  :=\;&f(\bx)+h(\by)
  +\langle\blmbd,\calA(\bx)-\by\rangle \\
  &+\frac{\rho}{2}\|\calA(\bx)-\by\|^2
  -\frac{1}{2s_k}\|\blmbd-\bmu\|^2.
  \end{aligned}
\end{equation*}

For fixed $(\bx,\bmu)$, $\mathcal P_{s_k}$ is $\rho$-strongly convex in
$\by$ and $1/s_k$-strongly concave in $\blmbd$. Hence
\begin{equation*}
  \red{{}\min_{\by}\max_{\blmbd}\mathcal P_{s_k}}
  =
  \green{\max_{\blmbd}\min_{\by}\mathcal P_{s_k}},
\end{equation*}
Both problems attain the same \blue{unique} saddle point
$\bigl(\by^\star,\blmbd^\star\bigr)$:
\hyperlink{appendix:why-rho-plus-sk}{%
    \beamergotobutton{Why $\rho+s_k$?}%
  }
\begin{equation*}
  \begin{aligned}
  \by^\star
  &=\prox_{h/(\rho+s_k)}
  \left(
    \calA(\bx)+\frac{\bmu}{\rho+s_k}
  \right), \\
  \blmbd^\star
  &=\bmu+s_k\bigl(\calA(\bx)-\by^\star\bigr).
  \end{aligned}
\end{equation*}

Two different one-step update strategies:

\begin{itemize}
  \item \red{DDLALM}: solve the saddle problem \red{\textbf{exactly.}} The answer seems to be \red{not!} 
 
  \item Candidate route: solve the saddle problem \green{\textbf{inexactly.}}

\end{itemize}


\end{frame}

\begin{frame}{Candidate Solution: Reduced-Dual Route}
\small
\begin{itemize}
  \item[\blue{$\triangleright$}]
  \textbf{Eliminate the $\by$-block through the $\rho$-prox:}
  \begin{equation*}
    d_s(\bx,\blmbd;\bmu)
    :=\min_{\by}\mathcal P_s(\bx,\by,\blmbd;\bmu),
    \qquad
    \red{\by(\blmbd;\bx)}
    =\operatorname{prox}_{h/\rho}
    \left(\calA(\bx)+\frac{\blmbd}{\rho}\right).
  \end{equation*}

  \item[\blue{$\triangleright$}]
  \textbf{Obtain an exact subdifferential certificate:}
  \begin{equation*}
    \blmbd+\rho\bigl[\calA(\bx)-\by(\blmbd;\bx)\bigr]
    \in\partial h\bigl(\by(\blmbd;\bx)\bigr).
  \end{equation*}

  \item[\blue{$\triangleright$}]
  \textbf{Apply Danskin's theorem:}
  \hfill
  \hyperlink{appendix:danskin-maxmin}{%
    \beamergotobutton{Danskin derivation $\rightarrow$}%
  }
  \begin{equation*}
    \nabla_{\bx}d_s
    =\nabla f(\bx)+J_{\calA}(\bx)^\top
    \left\{\blmbd+\rho\bigl[\calA(\bx)-\by(\blmbd;\bx)\bigr]\right\},
  \end{equation*}
  \begin{equation*}
    \nabla_{\blmbd}d_s
    =\calA(\bx)-\by(\blmbd;\bx)-\frac{\blmbd-\bmu}{s},
    \qquad
    \nabla_{\bmu}d_s
    =\frac{\blmbd-\bmu}{s}.
  \end{equation*}
\end{itemize}

\begin{block}{Moving-center identity}
\begin{equation*}
  \blue{\boxed{
  \calA(\bx)-\by(\blmbd;\bx)
  =\nabla_{\blmbd}d_s(\bx,\blmbd;\bmu)
  +\nabla_{\bmu}d_s(\bx,\blmbd;\bmu)}}
\end{equation*}
The $\bmu$-gradient exactly complements the regularization term in the
$\blmbd$-gradient.
\end{block}

\end{frame}

\begin{frame}{Candidate Algorithm}
\footnotesize

\noindent\textbf{Idea:} Apply $\bx$-descent and $(\blmbd,\bmu)$-ascent steps to the value function $d_{s_k}$.

\noindent\textbf{Parameters:} tolerance $\varepsilon>0$, stepsizes
$\eta_k,\alpha_k,\tau_k$, regularization $s_k$, and maximum iteration count $K$.

\vspace{0.08cm}
\begin{itemize}
  \item[\blue{$\triangleright$}]
  \textbf{Use the same $\rho$-prox mapping throughout:}
  \begin{equation*}
    \by(\blmbd;\bx)
    =\operatorname{prox}_{h/\rho}
    \left(\calA(\bx)+\frac{\blmbd}{\rho}\right).
  \end{equation*}

  \item[\blue{$\triangleright$}]
  \textbf{$\bx$-gradient descent on $d_{s_k}(\cdot,\blmbd^k;\bmu^k)$:}
  \begin{equation*}
    \bx^{k+1}=\bx^k-\eta_k
    \left\{\nabla f(\bx^k)+J_{\calA}(\bx^k)^\top
    \left[\blmbd^k+\rho\bigl(\calA(\bx^k)-\by(\blmbd^k;\bx^k)\bigr)\right]\right\}.
  \end{equation*}

  \item[\blue{$\triangleright$}]
  \textbf{One reduced-dual $\blmbd$-gradient ascent step:}
  \begin{equation*}
    \blmbd^{k+1}=\blmbd^k+\alpha_k
    \left[\calA(\bx^{k+1})-\by(\blmbd^k;\bx^{k+1})
    -\frac{\blmbd^k-\bmu^k}{s_k}\right].
  \end{equation*}

  \item[\blue{$\triangleright$}]
  \textbf{Moving-center $\bmu$-gradient ascent:}
  \begin{equation*}
    \bmu^{k+1}
    =\bmu^k+\tau_k\nabla_{\bmu}
    d_{s_k}(\bx^{k+1},\blmbd^{k+1};\bmu^k)
    =\bmu^k+\frac{\tau_k}{s_k}(\blmbd^{k+1}-\bmu^k),
    \qquad 0<\tau_k<s_k.
  \end{equation*}
\end{itemize}

\begin{block}{Oracle cost}
The $\bx$-first order evaluates the same $\rho$-prox mapping at
$(\blmbd^k,\bx^k)$ and $(\blmbd^k,\bx^{k+1})$.
\end{block}
\end{frame}

\begin{frame}{Exact $\blmbd$-Block Solve as a Special Case}
\small
\begin{itemize}
  \item[\blue{$\triangleright$}]
  \textbf{Choose the reduced-dual stepsize $\alpha_k=s_k$:}
  \begin{equation*}
    \begin{aligned}
    \blmbd^{k+1}
    &=\blmbd^k+s_k
    \left[\calA(\bx^{k+1})-\by(\blmbd^k;\bx^{k+1})
    -\frac{\blmbd^k-\bmu^k}{s_k}\right] \\
    &=\bmu^k+s_k
    \bigl[\calA(\bx^{k+1})-\by(\blmbd^k;\bx^{k+1})\bigr].
    \end{aligned}
  \end{equation*}

  \item[\blue{$\triangleright$}]
  \textbf{This is an exact solve of the frozen-$\by$ multiplier block:}
  \begin{equation*}
    \blmbd^{k+1}
    =\arg\max_{\blmbd}\,
    \mathcal P_{s_k}
    \bigl(\bx^{k+1},\by(\blmbd^k;\bx^{k+1}),\blmbd;\bmu^k\bigr).
  \end{equation*}

  \item[\blue{$\triangleright$}]
  \textbf{The distinction from reduced-dual exact maximization:}
  in $d_{s_k}(\bx^{k+1},\blmbd;\bmu^k)$, the minimizer
  $\by(\blmbd;\bx^{k+1})$ changes with $\blmbd$.
\end{itemize}

\begin{block}{Interpretation}
$\alpha_k=s_k$ is exact for the frozen-$\by$ block, but it is not an exact
maximizer of $d_{s_k}$.
\end{block}
\end{frame}

\begin{frame}{The role of $\alpha_k$}
\footnotesize
\begin{itemize}
  \item[\blue{$\triangleright$}]
  \textbf{Core bottleneck:}

 \(\text{primal descent}
\le
-\text{descent}
+
C\eta_k\|\lambda^k-\lambda_k^\star(x^k)\|^2
+
\text{regularization drift},\) which limits the $\eta_k$.
  \item[\blue{$\triangleright$}]
  \textbf{The frozen-$\by$ exact step is not the stable reduced-dual scale:}
  \begin{equation*}
   L_{\lambda,k}
\leq
\frac1\rho+\frac1{s_k},
\qquad
\alpha_k
\leq
\frac{1}{
\frac1\rho+\frac1{s_k}
}
=
\frac{\rho s_k}{\rho+s_k}
\leq
\frac1{L_{\lambda,k}}.
  \end{equation*}
  Thus $\alpha_k=s_k$ does not respect this scale when $s_k\uparrow$.

  \item[\blue{$\triangleright$}]
  \textbf{Proposed stable dual step:}
  \begin{equation*}
    \blue{\alpha_k=\theta_\lambda\frac{\rho s_k}{\rho+s_k}},
    \qquad 0<\theta_\lambda\leq1.
  \end{equation*}
  It yields the convex-combination representation
  \begin{equation*}
    \begin{aligned}
    \blmbd^{k+1}
    ={}&(1-\theta_\lambda)\blmbd^k
    +\theta_\lambda\frac{s_k}{\rho+s_k}
    \left\{\blmbd^k+\rho\bigl[\calA(\bx^{k+1})
    -\by(\blmbd^k;\bx^{k+1})\bigr]\right\} \\
    &+\theta_\lambda\frac{\rho}{\rho+s_k}\bmu^k.
    \end{aligned}
  \end{equation*}

  \item[\blue{$\triangleright$}]
  \textbf{Main parameter schedule:}
  \begin{equation*}
    s_k=c_s(k+1)^a,\qquad
    \eta_k=\frac{\theta_x}{\rho+s_k},\qquad
    \alpha_k=\theta_\lambda\frac{\rho s_k}{\rho+s_k},\qquad
    \tau_k=\theta_\mu s_k,
  \end{equation*}
  where $a\in(0,1)$, $c_s>0$, $0<\theta_\lambda\leq1$, and
  $0<\theta_\mu<1$. The complexity target uses $a=1/3$.
\end{itemize}

\begin{block}{Question}
Can stable dual contraction and constant-rate center tracking jointly absorb
the multiplier tracking error?
\end{block}
\end{frame}

% \begin{frame}{Proof Roadmap: Four Coupled Modules}
% \small
% \begin{itemize}
%   \item[\blue{$\triangleright$}]
%   \textbf{Module 1: Basic identities and multiplier boundedness.}
%   Use the exact $\rho$-prox certificate and the convex-combination updates to
%   bound $\{\blmbd^k\}$ and $\{\bmu^k\}$ uniformly.

%   \item[\blue{$\triangleright$}]
%   \textbf{Module 2: Primal descent and multiplier tracking.}
%   Define the exact regularized target
%   $\blmbd_{s_k}^{\star}(\bx;\bmu):=\arg\max_{\blmbd}d_{s_k}(\bx,\blmbd;\bmu)$
%   and derive a recursion for
%   \begin{equation*}
%     e_k:=\blmbd^k-\blmbd_{s_k}^{\star}(\bx^k;\bmu^k).
%   \end{equation*}

%   \item[\blue{$\triangleright$}]
%   \textbf{Module 3: Center-gap control and a joint Lyapunov function.}
%   Track $\delta_k:=\blmbd^k-\bmu^k$ and use its decrease to control the
%   moving-center drift and the regularization drift.

%   \item[\blue{$\triangleright$}]
%   \textbf{Module 4: KKT reconstruction and complexity.}
%   Combine the primal step with
%   \begin{equation*}
%     \calA(\bx)-\by(\blmbd;\bx)
%     =\nabla_{\blmbd}d_s+\nabla_{\bmu}d_s
%   \end{equation*}
%   to recover stationarity and feasibility at a tail-best iterate.
% \end{itemize}
% \end{frame}

% \begin{frame}{Can the Moving Center Absorb the Tracking Error?}
% \small
% \begin{itemize}
%   \item[\blue{$\triangleright$}]
%   \textbf{Two coupled errors must be controlled:}
%   \begin{equation*}
%     e_k=\blmbd^k-\blmbd_{s_k}^{\star}(\bx^k;\bmu^k),
%     \qquad
%     \delta_k=\blmbd^k-\bmu^k.
%   \end{equation*}

%   \item[\blue{$\triangleright$}]
%   \textbf{Candidate joint Lyapunov function:}
%   \begin{equation*}
%     \mathcal V_k
%     =\psi_{s_k}(\bx^k;\bmu^k)
%     +c_e\|e_k\|^2+c_\delta\|\delta_k\|^2.
%   \end{equation*}

%   \item[\blue{$\triangleright$}]
%   \textbf{Decisive inequality to prove:}
%   find positive coefficients such that
%   \begin{equation*}
%     \mathcal V_{k+1}-\mathcal V_k
%     \leq-c_x\|\Delta_{\bx}^{k+1}\|^2
%     -c_e'\|e_k\|^2-c_\delta'\|\delta_k\|^2
%     +\mathcal E_k^s,
%   \end{equation*}
%   where $\mathcal E_k^s$ is the controllable drift induced by
%   $s_{k+1}-s_k$.
% \end{itemize}

% \begin{block}{Hu--Jiang open problem}
% Can the center-gap decrease absorb multiplier tracking and moving-center drift
% simultaneously while retaining a single loop?
% \end{block}
% \end{frame}


% \section{References}

\begin{frame}{References}
\footnotesize
\begin{enumerate}
  \item[{[1]}] H. Zhang, N. Marumo, T. K. Pong, and A. Takeda,
  ``Complexity and Convergence Analysis of a Single-Loop SDCAM for
  Lipschitz Composite Optimization and Beyond,'' arXiv:2512.24059, 2025.

  \item[{[2]}] K. Sun and X. A. Sun,
  ``Dual Descent Augmented Lagrangian Method and Alternating Direction
  Method of Multipliers,'' \emph{SIAM Journal on Optimization},
  34(2):1679--1707, 2024.

  \item[{[3]}] J. M. Danskin,
  ``The Theory of Max--Min, with Applications,''
  \emph{SIAM Journal on Applied Mathematics},
  14(4):641--664, 1966.
\end{enumerate}
\end{frame}

\appendix

\begin{frame}{Appendix Contents}
\hypertarget{appendix:contents}{}
\small
\begin{enumerate}
  \setlength{\itemsep}{0.22cm}
  \item \hyperlink{appendix:lemmas}{\textbf{Appendix I:} Lemmas 1--4}
  \item \hyperlink{appendix:bootstrap-joint-claim}{\textbf{Appendix II:} Induction}
  \item \hyperlink{appendix:danskin-theorem}{\textbf{Appendix III:} Danskin's Theorem}
  \item \hyperlink{appendix:why-rho-plus-sk}{\textbf{Appendix IV:} Why $\rho+s_k$ Appears}
  \item \hyperlink{appendix:danskin-maxmin}{\textbf{Appendix V:} Danskin's Theorem for the Max--Min Problem}
  \item \hyperlink{appendix:exact-maxmin}{\textbf{Appendix VI:} Exact Solution of the Max--Min Problem}
\end{enumerate}
\end{frame}

% BEGIN APPENDIX I LEMMAS
\begin{frame}{Appendix I: Lemma 1---Multiplier and Feasibility Bounds
\hfill\hyperlink{appendix:contents}{\beamerreturnbutton{Contents}}}
\hypertarget{appendix:lemmas}{}
\label{app:lemmas}
\scriptsize
Suppose $0<\tau_k\le s_k$. Define
\begin{equation*}
  M_\mu:=\max\{\|\bmu^0\|,L_h\},
  \qquad C_r:=L_h+M_\mu.
\end{equation*}

\noindent\textbf{Conclusion:} For every $k\ge0$,
\begin{align}
  \|\blmbd^{k+1}\|\le M_\mu,
  \qquad
  \|\bmu^{k+1}\|\le M_\mu,
  \label{eq:two-multiplier-bounds}\\
  \|\br^{k+1}\|
  &\le\frac{C_r}{\rho_{s_k}},
  \qquad
  \frac{\|\bmu^{k+1}-\bmu^k\|^2}{\tau_k}
  \le\frac{C_r^2\tau_k}{\rho_{s_k}^2}.
  \label{eq:feasibility-drift-control}
\end{align}

\noindent\textbf{Proof:}
\begin{itemize}
\setlength{\itemsep}{0.06cm}

\item[\blue{$\triangleright$}]
\textbf{Saddle certificate.}
Assumption (iii) and the saddle certificate in
\eqref{eq:key-exact-identities} give
\begin{equation*}
  \blmbd^{k+1}+\rho\br^{k+1}\in\partial h(\by^{k+1}),
  \qquad
  \|\blmbd^{k+1}+\rho\br^{k+1}\|\le L_h.
\end{equation*}

\item[\blue{$\triangleright$}]
\textbf{Two-multiplier boundedness.}
The base case holds since $\|\bmu^0\|\le M_\mu$.
Assume $\|\bmu^k\|\le M_\mu$.
Applying the two convex combinations in
\eqref{eq:key-convex-combinations}, respectively, gives
\begin{equation*}
\begin{aligned}
  \|\blmbd^{k+1}\|
  &\le \frac{\rho}{\rho_{s_k}}M_\mu
  +\frac{s_k}{\rho_{s_k}}L_h
  \le M_\mu, \\
  \|\bmu^{k+1}\|
  &\le \left(1-\frac{\tau_k}{s_k}\right)M_\mu
  +\frac{\tau_k}{s_k}M_\mu
  =M_\mu.
\end{aligned}
\end{equation*}
Thus, \eqref{eq:two-multiplier-bounds} follows by induction.

\item[\blue{$\triangleright$}]
\textbf{Feasibility and center-multiplier drift.}
Since $\rho_{s_k}\br^{k+1}
 =(\blmbd^{k+1}+\rho\br^{k+1})-\bmu^k$ and
$\bmu^{k+1}-\bmu^k=\tau_k\br^{k+1}$ by
\eqref{eq:key-exact-identities},
\begin{equation*}
\begin{aligned}
  \|\br^{k+1}\|
  &\le\frac{L_h+M_\mu}{\rho_{s_k}}
  =\frac{C_r}{\rho_{s_k}},\\
  \frac{\|\bmu^{k+1}-\bmu^k\|^2}{\tau_k}
  &=\tau_k\|\br^{k+1}\|^2
  \leq\frac{C_r^2\tau_k}{\rho_{s_k}^2}.
\end{aligned}
\end{equation*}
\end{itemize}

\end{frame}

\begin{frame}{Appendix I: Lemma 2---Quasi-Descent on $\Phi_s$
\hfill\hyperlink{appendix:contents}{\beamerreturnbutton{Contents}}}
\vspace{0.2cm}
\scriptsize


Let $\Delta_{\bx}^{k+1}:=\bx^{k+1}-\bx^k$, and define
$\Delta_{\by}^{k+1}$ and $\Delta_{\bmu}^{k+1}$ similarly.

\noindent\textbf{Conclusion:}
\begin{equation}\label{eq:slide-quasi-descent}
\begin{split}
  &\Phi_{s_{k+1}}(\bx^{k+1},\by^{k+1};\bmu^{k+1})
  -\Phi_{s_k}(\bx^k,\by^k;\bmu^k)\\
  &\quad\leq
  -\frac{3\rho_{s_k}}{4\theta_x}\|\Delta_{\bx}^{k+1}\|^2
  -\frac{\rho_{s_k}}{2}\|\Delta_{\by}^{k+1}\|^2
  +\left(\frac{1}{\tau_k}
  +\frac{\rho_{s_{k+1}}-\rho_{s_k}}{2\tau_k^2}\right)
  \|\Delta_{\bmu}^{k+1}\|^2.
\end{split}
\end{equation}

\noindent\textbf{Proof:}

\begin{itemize}
\setlength{\itemsep}{0.08cm}
\item[\blue{$\triangleright$}]
\textbf{$\bx$-block descent.}
With $\eta_k=\theta_x/\rho_{s_k}$, the $\bx$-gradient step gives
\begin{equation*}
  \Phi_{s_k}(\bx^{k+1},\by^k;\bmu^k)
  -\Phi_{s_k}(\bx^k,\by^k;\bmu^k)
  \leq
  -\frac{3\rho_{s_k}}{4\theta_x}
  \|\Delta_{\bx}^{k+1}\|^2.
\end{equation*}

\item[\blue{$\triangleright$}]
\textbf{$\by$-block descent.}
The exact $(\by,\blmbd)$ saddle solve makes $\by^{k+1}$ the exact minimizer of the strongly convex $\by$-subproblem. Hence,
\begin{equation*}
  \Phi_{s_k}(\bx^{k+1},\by^{k+1};\bmu^k)
  -\Phi_{s_k}(\bx^{k+1},\by^k;\bmu^k)
  \leq
  -\frac{\rho_{s_k}}{2}
  \|\Delta_{\by}^{k+1}\|^2.
\end{equation*}

\item[\blue{$\triangleright$}]
\textbf{Center-multiplier shift.}
Using $\Delta_{\bmu}^{k+1}=\tau_k\br^{k+1}$ from
\eqref{eq:key-exact-identities},
\begin{equation*}
  \Phi_{s_k}(\bx^{k+1},\by^{k+1};\bmu^{k+1})
  -\Phi_{s_k}(\bx^{k+1},\by^{k+1};\bmu^k)
  =\frac{1}{\tau_k}\|\Delta_{\bmu}^{k+1}\|^2.
\end{equation*}

\item[\blue{$\triangleright$}]
\textbf{Regularization shift.}
The variation of $s_k$ contributes
\begin{equation*}
  \Phi_{s_{k+1}}(\bx^{k+1},\by^{k+1};\bmu^{k+1})
  -\Phi_{s_k}(\bx^{k+1},\by^{k+1};\bmu^{k+1})
  =\frac{\rho_{s_{k+1}}-\rho_{s_k}}{2\tau_k^2}
  \|\Delta_{\bmu}^{k+1}\|^2.
\end{equation*}
\end{itemize}

\end{frame}

\begin{frame}{Appendix I: Lemma 3---KKT Residual Control
\hfill\hyperlink{appendix:contents}{\beamerreturnbutton{Contents}}}
\scriptsize

Fix $k\in\{\lceil K/2\rceil,\ldots,K-1\}$. Then
\begin{equation*}
\begin{aligned}
  &\blmbd^k+\rho\br^k\in\partial h(\by^k),\\
  &\left\|\nabla f(\bx^k)+J_{\calA}(\bx^k)^\top
  \bigl(\blmbd^k+\rho\br^k\bigr)\right\|
  \le \eta_k^{-1}\|\Delta_{\bx}^{k+1}\|
  +G_{\calA}C_r\left(
  \frac{\tau_{k-1}}{\rho_{s_{k-1}}}
  +\frac{\rho_{s_k}-\rho_{s_{k-1}}}{\rho_{s_{k-1}}}
  \right),\\
  &\|\calA(\bx^k)-\by^k\| \le\frac{C_r}{\rho_{s_{k-1}}}
\end{aligned}
\end{equation*}

\noindent\textbf{Proof:}
\begin{itemize}
\setlength{\itemsep}{0.06cm}
\item[\blue{$\triangleright$}]
\textbf{Choose the KKT multiplier.}
By \eqref{eq:key-exact-identities} at iteration $k-1$,
\begin{equation}\label{eq:kkt-multiplier}
  \blmbd^k+\rho\br^k
  =\bmu^{k-1}+\rho_{s_{k-1}}\br^k
  \in\partial h(\by^k).
\end{equation}

\item[\blue{$\triangleright$}]
\textbf{$\bx$-step optimality.}
The $\bx$-update gives
\begin{equation}\label{eq:kkt-x-optimality}
\begin{aligned}
  \zero
  &=\eta_k^{-1}\Delta_{\bx}^{k+1}
  +\nabla_{\bx}\Phi_{s_k}
  (\bx^k,\by^k;\bmu^k)\\
  &=\eta_k^{-1}\Delta_{\bx}^{k+1}
  +\nabla f(\bx^k)
  +J_{\calA}(\bx^k)^\top
  \bigl(\bmu^k+\rho_{s_k}\br^k\bigr).
\end{aligned}
\end{equation}

\item[\blue{$\triangleright$}]
\textbf{Multiplier mismatch.}
Using $\bmu^k-\bmu^{k-1}=\tau_{k-1}\br^k$,
\begin{equation*}
\begin{aligned}
  &\bmu^k+\rho_{s_k}\br^k
  -\bigl(\blmbd^k+\rho\br^k\bigr)\\
  &\quad=\bigl(\tau_{k-1}
  +\rho_{s_k}-\rho_{s_{k-1}}\bigr)\br^k.
\end{aligned}
\end{equation*}
Hence, \eqref{eq:feasibility-drift-control} and
$\|J_{\calA}(\bx^k)\|\le G_{\calA}$ give the results.
% \begin{equation*}
% \begin{aligned}
%   &\left\|\nabla f(\bx^k)+J_{\calA}(\bx^k)^\top
%   \bigl(\blmbd^k+\rho\br^k\bigr)\right\|\\
%   &\quad\le \eta_k^{-1}\|\Delta_{\bx}^{k+1}\|
%   +G_{\calA}C_r\left(
%   \frac{\tau_{k-1}}{\rho_{s_{k-1}}}
%   +\frac{\rho_{s_k}-\rho_{s_{k-1}}}{\rho_{s_{k-1}}}
%   \right)\\
%   &\quad=\eta_k^{-1}\|\Delta_{\bx}^{k+1}\|+O(K^{-1}),
%   \qquad
%   \|\br^k\|\le\frac{C_r}{\rho_{s_{k-1}}}=O(K^{-a}).
% \end{aligned}
% \end{equation*}
\end{itemize}
\end{frame}


\begin{frame}{Appendix I: Lemma 4---Weighted Finite Length
\hfill\hyperlink{appendix:contents}{\beamerreturnbutton{Contents}}}
\scriptsize

Let
$D_k:=\|\Delta_{\bx}^{k+1}\|^2+\|\Delta_{\by}^{k+1}\|^2$ and
$c_{\mathrm{desc}}:=\min\{3/(4\theta_x),1/2\}>0$.

\noindent\textbf{Condition:} Let $s_k=c_s(k+1)^a$ and
$\tau_k=\theta_\mu s_k/(k+2)$ with $a\in(0,1)$. Then
\begin{equation}\label{eq:slide-weighted-finite-length}
  c_{\mathrm{desc}}\sum_{k=0}^{K-1}
  \rho_{s_k}\|\Delta_{\bx}^{k+1}\|^2
  \le c_{\mathrm{desc}}\sum_{k=0}^{K-1}\rho_{s_k}D_k
  \le C_{\mathrm{desc}}<\infty.
\end{equation}

\noindent\textbf{Proof:}
\begin{itemize}
\setlength{\itemsep}{0.06cm}
\item[\blue{$\triangleright$}]
\textbf{The moving-potential error is summable.}
\begin{equation*}
\begin{aligned}
  e_k
  &:=C_r^2\left(
  \frac{\tau_k}{\rho_{s_k}^2}
  +\frac{\rho_{s_{k+1}}-\rho_{s_k}}{2\rho_{s_k}^2}
  \right),\\
  \frac{\tau_k}{\rho_{s_k}^2}
  &=O(k^{-(1+a)}),
  \qquad
  \frac{\rho_{s_{k+1}}-\rho_{s_k}}{\rho_{s_k}^2}
  =O(k^{-(1+a)}).
\end{aligned}
\end{equation*}
Hence, $\sum_{k=0}^{\infty}e_k<\infty$.

\item[\blue{$\triangleright$}]
\textbf{Quasi-descent leaves a finite error budget.}
Lemma 2 and Lemma 1 yield
\begin{equation}\label{eq:slide-clean-quasi-descent}
  \Phi_{s_{k+1}}^{k+1}-\Phi_{s_k}^{k}
  \le -c_{\mathrm{desc}}\rho_{s_k}D_k+e_k,
  \qquad
  \Phi_{s_k}^{k}:=\Phi_{s_k}(\bx^k,\by^k;\bmu^k).
\end{equation}

\item[\blue{$\triangleright$}]
\textbf{Telescoping converts descent into a global budget.}
Summing \eqref{eq:slide-clean-quasi-descent} and using 
$\Phi_{s_K}^{K}\ge\underline\Phi$ give
\begin{equation*}
\begin{aligned}
  c_{\mathrm{desc}}\sum_{k=0}^{K-1}\rho_{s_k}D_k
  % &\le \Phi_{s_0}^{0}-\Phi_{s_K}^{K}
  % +\sum_{k=0}^{K-1}e_k\\
  \le \Phi_{s_0}^{0}-\underline\Phi
  +\sum_{k=0}^{\infty}e_k
  =:C_{\mathrm{desc}}<\infty.
\end{aligned}
\end{equation*}
\end{itemize}
\end{frame}

% END APPENDIX I LEMMAS

% \section{Appendix}

\begin{frame}{Appendix II.1: Induction---Joint Claim
\hfill\hyperlink{appendix:contents}{\beamerreturnbutton{Contents}}}
\hypertarget{appendix:bootstrap-joint-claim}{}
\scriptsize
\begin{itemize}
\setlength{\itemsep}{0.11cm}

\item[\blue{$\triangleright$}]
\textbf{Summable drift budget.}
Define
\begin{equation*}
  e_k:=C_r^2\left(
  \frac{\tau_k}{\rho_{s_k}^2}
  +\frac{\rho_{s_{k+1}}-\rho_{s_k}}{2\rho_{s_k}^2}
  \right),
  \qquad
  C_e:=\sum_{k=0}^{\infty}e_k<\infty.
\end{equation*}

\item[\blue{$\triangleright$}]
\textbf{Derived level set.}
Set
\begin{equation*}
  \overline\Phi:=\Phi_{s_0}(\bx^0,\by^0;\bmu^0)+C_e,
  \qquad
  \alpha:=\overline\Phi+\frac{C_r^2}{2\rho_{s_0}},
  \qquad
  \calS_\alpha:=\{\bx:F(\bx)\leq\alpha\}.
\end{equation*}
Choose a compact convex set $\calS$ such that
$\calS_\alpha\subseteq\operatorname{int}(\calS)$.

\item[\blue{$\triangleright$}]
\textbf{Potential-to-objective bridge.}
Multiplier boundedness and the Lipschitz continuity of $h$ give
\begin{equation*}
  F(\bx^k)
  \leq
  \Phi_{s_k}(\bx^k,\by^k;\bmu^k)
  +\frac{C_r^2}{2\rho_{s_k}}.
\end{equation*}

\item[\blue{$\triangleright$}]
\textbf{Joint induction claim.}
For every $k\geq0$, prove simultaneously that
\begin{equation*}
  \Phi_{s_k}(\bx^k,\by^k;\bmu^k)
  \leq
  \Phi_{s_0}(\bx^0,\by^0;\bmu^0)+\sum_{j=0}^{k-1}e_j,
  \qquad
  \bx^k\in\calS_\alpha.
\end{equation*}

\item[\blue{$\triangleright$}]
\textbf{Base case.}
The potential bound is immediate, and
\begin{equation*}
  F(\bx^0)
  \leq\Phi_{s_0}(\bx^0,\by^0;\bmu^0)
  +\frac{C_r^2}{2\rho_{s_0}}
  \leq\alpha.
\end{equation*}
Hence, $\bx^0\in\calS_\alpha$.
\end{itemize}
\end{frame}

\begin{frame}{Appendix II.2: Induction---Smooth Region
\hfill\hyperlink{appendix:contents}{\beamerreturnbutton{Contents}}}
\scriptsize
\begin{itemize}
\setlength{\itemsep}{0.13cm}

\item[\blue{$\triangleright$}]
\textbf{Induction hypothesis.}
Assume that the joint claim holds at iteration $k$. In particular,
\begin{equation*}
  \bx^k\in\calS_\alpha\subseteq\operatorname{int}(\calS).
\end{equation*}

\item[\blue{$\triangleright$}]
\textbf{Safety margin.}
Define
\begin{equation*}
  d_{\calS}:=\operatorname{dist}(\calS_\alpha,\partial\calS)>0,
  \qquad
  M_d:=\max\left\{
  \frac{\|\nabla_{\bx}\Phi_{s_0}(\bx^0,\by^0;\bmu^0)\|}{\rho_{s_0}},
  \frac{G_f+G_{\calA}(M_\mu+C_r)}{\rho_{s_0}}
  \right\}.
\end{equation*}

\item[\blue{$\triangleright$}]
\textbf{Normalized gradient bound.}
Lemma 1 and Assumption (ii) imply
\begin{equation*}
  \frac{\|\nabla_{\bx}\Phi_{s_k}(\bx^k,\by^k;\bmu^k)\|}
       {\rho_{s_k}}
  \leq M_d.
\end{equation*}
The definition of $M_d$ covers $k=0$; the bounds
$\|\bmu^k\|\leq M_\mu$ and $\|\br^k\|\leq C_r/\rho_{s_{k-1}}$
cover $k\geq1$.

\item[\blue{$\triangleright$}]
\textbf{The next iterate remains in $\calS$.}
Choose $\eta_k=\theta_x/\rho_{s_k}$ with
$\theta_xM_d\leq d_{\calS}$. Then
\begin{equation*}
  \|\bx^{k+1}-\bx^k\|
  =\frac{\theta_x}{\rho_{s_k}}
  \|\nabla_{\bx}\Phi_{s_k}(\bx^k,\by^k;\bmu^k)\|
  \leq\theta_xM_d
  \leq d_{\calS}.
\end{equation*}
Thus, $\bx^{k+1}\in\calS$. Convexity also gives
$[\bx^k,\bx^{k+1}]\subseteq\calS$.

\item[\blue{$\triangleright$}]
\textbf{Primal descent is now valid.}
The whole line segment lies in the smoothness region. Therefore, the
one-step descent inequality of Lemma 2 applies at iteration $k$.
\end{itemize}
\end{frame}

\begin{frame}{Appendix II.3: Induction---Closure
\hfill\hyperlink{appendix:contents}{\beamerreturnbutton{Contents}}}
\scriptsize
\begin{itemize}
\setlength{\itemsep}{0.16cm}

\item[\blue{$\triangleright$}]
\textbf{Apply quasi-descent.}
Since $\bx^k,\bx^{k+1}\in\calS$, Lemma 2 and the drift bound give
\begin{equation*}
\begin{aligned}
  \Phi_{s_{k+1}}(\bx^{k+1},\by^{k+1};\bmu^{k+1})
  &\leq \Phi_{s_k}(\bx^k,\by^k;\bmu^k)+e_k\\
  &\leq \Phi_{s_0}(\bx^0,\by^0;\bmu^0)
  +\sum_{j=0}^{k}e_j.
\end{aligned}
\end{equation*}

\item[\blue{$\triangleright$}]
\textbf{Convert potential control to objective control.}
Using $\sum_{j=0}^{k}e_j\leq C_e$ and
$\rho_{s_{k+1}}\geq\rho_{s_0}$,
\begin{equation*}
\begin{aligned}
  F(\bx^{k+1})
  &\leq
  \Phi_{s_{k+1}}(\bx^{k+1},\by^{k+1};\bmu^{k+1})
  +\frac{C_r^2}{2\rho_{s_{k+1}}}\\
  &\leq
  \overline\Phi+\frac{C_r^2}{2\rho_{s_0}}
  =\alpha.
\end{aligned}
\end{equation*}

\item[\blue{$\triangleright$}]
\textbf{Return to the derived level set.}
The definition of $\calS_\alpha$ gives
\begin{equation*}
  \bx^{k+1}\in\calS_\alpha.
\end{equation*}
Both parts of the joint claim now hold at iteration $k+1$.
\end{itemize}

\begin{block}{mechanism}
\centering
$\displaystyle
\bx^k\in\calS_\alpha
\ \Longrightarrow\ 
\bx^{k+1}\in\calS
\ \Longrightarrow\ 
\text{quasi-descent is valid}
\ \Longrightarrow\ 
F(\bx^{k+1})\leq\alpha
\ \Longrightarrow\ 
\bx^{k+1}\in\calS_\alpha.$
\end{block}
\end{frame}

\begin{frame}{Appendix III: Danskin's Theorem
\hfill\hyperlink{appendix:contents}{\beamerreturnbutton{Contents}}}
\hypertarget{appendix:danskin-theorem}{}
\small
Let
\begin{equation*}
  V(\mathbf u):=\max_{\mathbf v\in\mathcal V}g(\mathbf u,\mathbf v),
  \qquad
  \mathbf v^\star(\mathbf u)
  \in\operatorname*{arg\,max}_{\mathbf v\in\mathcal V}
  g(\mathbf u,\mathbf v).
\end{equation*}
Assume that the maximum is attained, $g$ is continuously differentiable in
$\mathbf u$, and the standard continuity conditions of Danskin's
theorem~{\normalfont[3]} hold.

\begin{block}{Unique-maximizer form}
If $\mathbf v^\star(\mathbf u)$ is unique, then $V$ is differentiable at
$\mathbf u$ and
\begin{equation*}
  \blue{\nabla V(\mathbf u)=\nabla_{\mathbf u}g
  \bigl(\mathbf u,\mathbf v^\star(\mathbf u)\bigr).}
\end{equation*}
No derivative of the solution map $\mathbf v^\star(\mathbf u)$ is required.
\end{block}

% \begin{block}{Application to $\Phi_s$}
% The $1/s$-strong concavity of
% $\blmbd\mapsto\mathcal P_s(\bx,\by,\blmbd;\bmu)$ gives the unique maximizer
% $\blmbd_s$. Hence, for $\mathbf u=\bx$ or $\mathbf u=\bmu$,
% \begin{equation*}
%   \nabla_{\mathbf u}\Phi_s(\bx,\by;\bmu)
%   =\nabla_{\mathbf u}\mathcal P_s
%   \bigl(\bx,\by,\blmbd_s(\bx,\by;\bmu);\bmu\bigr).
% \end{equation*}
% \end{block}
\end{frame}

% Original derivation
\begin{frame}{Appendix IV: Why $\rho+s_k$ Appears
\hfill\hyperlink{appendix:contents}{\beamerreturnbutton{Contents}}}
\hypertarget{appendix:why-rho-plus-sk}{}

For fixed $(\bx^{k+1},\bmu^k)$, the two conditions are
\begin{equation*}
  0\in\partial h(\by^{k+1})-\blmbd^{k+1}
  +\rho\bigl(\by^{k+1}-\calA(\bx^{k+1})\bigr),
\end{equation*}
\begin{equation*}
  \blmbd^{k+1}=\bmu^k+s_k
  \bigl(\calA(\bx^{k+1})-\by^{k+1}\bigr).
\end{equation*}
Substituting the second condition into the first gives
\begin{equation*}
  0\in\partial h(\by^{k+1})-\bmu^k
  +(\rho+s_k)\bigl(\by^{k+1}-\calA(\bx^{k+1})\bigr),
\end{equation*}
and therefore the explicit $(\rho+s_k)$-prox update follows.
\end{frame}

\begin{frame}{Appendix V: Danskin's Theorem for the Max--Min Problem
\hfill\hyperlink{appendix:contents}{\beamerreturnbutton{Contents}}}
\small
For fixed $(\bx,\bmu)$, define the reduced dual function
\begin{equation*}
  d_s(\bx,\blmbd;\bmu)
  :=\min_{\by}\mathcal P_s(\bx,\by,\blmbd;\bmu).
\end{equation*}
The inner problem is $\rho$-strongly convex and has the unique minimizer
\begin{equation*}
  \by(\blmbd;\bx)
  =\operatorname{prox}_{h/\rho}
  \left(\calA(\bx)+\frac{\blmbd}{\rho}\right).
\end{equation*}
Thus, the max--min problem is equivalently
\begin{equation*}
  \max_{\blmbd}\min_{\by}\mathcal P_s(\bx,\by,\blmbd;\bmu)
  =\max_{\blmbd}d_s(\bx,\blmbd;\bmu).
\end{equation*}

\hypertarget{appendix:danskin-maxmin}{}
\begin{block}{Danskin's theorem}
Since the inner minimizer is unique and $\mathcal P_s$ is differentiable
with respect to $(\bx,\blmbd)$, Danskin's theorem gives
\begin{equation*}
  \begin{aligned}
  \nabla_{\blmbd}d_s(\bx,\blmbd;\bmu)
  &=\nabla_{\blmbd}\mathcal P_s
    \bigl(\bx,\by(\blmbd;\bx),\blmbd;\bmu\bigr) \\
  &=\calA(\bx)-\by(\blmbd;\bx)-\frac{1}{s}(\blmbd-\bmu).
  \end{aligned}
\end{equation*}
Moreover,
\begin{equation*}
  \begin{aligned}
  \nabla_{\bx}d_s(\bx,\blmbd;\bmu)
  &=\nabla_{\bx}\mathcal P_s
    \bigl(\bx,\by(\blmbd;\bx),\blmbd;\bmu\bigr) \\
  &=\nabla f(\bx)+J_{\calA}(\bx)^\top
    \left[\blmbd+\rho\bigl(\calA(\bx)-\by(\blmbd;\bx)\bigr)\right].
  \end{aligned}
\end{equation*}
No derivative of the possibly nonsmooth solution map $\by(\blmbd;\bx)$
with respect to either $\bx$ or $\blmbd$ is needed.
\end{block}
\end{frame}

\begin{frame}{Appendix VI: Exact Solution of the Max--Min Problem
\hfill\hyperlink{appendix:contents}{\beamerreturnbutton{Contents}}}
\hypertarget{appendix:exact-maxmin}{}
\small

The function $d_s(\bx,\cdot;\bmu)$ is $1/s$-strongly concave. Its unique
maximizer $\blmbd^\star$ therefore satisfies
\begin{equation*}
  0=\nabla_{\blmbd}d_s(\bx,\blmbd^\star;\bmu)
  =\calA(\bx)-\by(\blmbd^\star;\bx)
  -\frac{1}{s}(\blmbd^\star-\bmu).
\end{equation*}
Let $\by^\star:=\by(\blmbd^\star;\bx)$. The exact optimality conditions are
\begin{equation*}
  0\in\partial h(\by^\star)-\blmbd^\star
  +\rho\bigl(\by^\star-\calA(\bx)\bigr),
\end{equation*}
\begin{equation*}
  \blmbd^\star
  =\bmu+s\bigl(\calA(\bx)-\by^\star\bigr).
\end{equation*}
Substitution gives
\begin{equation*}
  0\in\partial h(\by^\star)-\bmu
  +(\rho+s)\bigl(\by^\star-\calA(\bx)\bigr).
\end{equation*}
Hence the unique exact solution is
\begin{equation*}
  \boxed{
  \by^\star=\operatorname{prox}_{h/(\rho+s)}
  \left(\calA(\bx)+\frac{\bmu}{\rho+s}\right),
  \qquad
  \blmbd^\star=\bmu+s\bigl(\calA(\bx)-\by^\star\bigr).}
\end{equation*}
% \begin{alertblock}{Consequence}
% Eliminating $y$ first exposes a $\rho$-prox, but exact maximization of the
% reduced dual function recovers the same $(\rho+s)$-dependent saddle point.
% \end{alertblock}
\end{frame}

\end{document}
